\section{Algebraic principalization and finite logarithmic charts}

\begin{lemma}[The affine blowup chart inside the fraction field]\label{lem:sc0-hironaka-affine-chart}
\lean{EnvelopeProperty.Char0.affineBlowUpReesChartEquiv,EnvelopeProperty.Char0.affineBlowUpReesChartEquiv_algebraMap}
\leanok
\uses{def:sc0-rees-chart-tree}
Let $C\subseteq K$ be a subalgebra of a field, let $g_1,\ldots,g_n\in C$, put
$I=(g_1,\ldots,g_n)$, and suppose $g_i\ne0$. The localization map $C[1/g_i]\to K$
induces an isomorphism
\[
 C[I/g_i]\simeq C[g_1/g_i,\ldots,g_n/g_i]\subseteq K
\]
which identifies the two maps from $C$.
\end{lemma}
\begin{proof}
The localization map sends $c/g_i^m$ to the same fraction in $K$ and is injective: if the
image is zero, multiplying by the nonzero denominator gives $c=0$. Its restriction to
$C[I/g_i]$ sends the generator $h/g_i$, with $h\in I$, into the right-hand algebra.
Indeed, write $h=\sum_j c_jg_j$ and divide this equality by $g_i$. Conversely each $g_j/g_i$
is the image of the corresponding generator of $C[I/g_i]$, and the image contains $C$.
Closure under sums and products proves equality of the image with the indicated algebra.
The identification of the maps from $C$ follows from the localization formula with exponent zero.
\end{proof}

\begin{lemma}[Principal covers as chart nodes]\label{lem:sc0-principal-cover-chart}
\lean{EnvelopeProperty.Char0.reesTreeChart_eq_adjoin_inv_of_span_eq_top,EnvelopeProperty.Char0.principalCoverLocalizationEquiv,EnvelopeProperty.Char0.principalCoverLocalizationEquiv_algebraMap}
\leanok
\uses{def:sc0-rees-chart-tree}
If $(h_1,\ldots,h_s)=C$, then for every nonzero $h_i$ one has
$C[h_1/h_i,\ldots,h_s/h_i]=C[1/h_i]$ inside $K$.
\end{lemma}
\begin{proof}
The ratios belong to $C[1/h_i]$. Choose $c_j\in C$ with $\sum_jc_jh_j=1$.
Then $1/h_i=\sum_jc_j(h_j/h_i)$ lies in the ratio algebra, proving the reverse inclusion.
Thus the chart node of this family is exactly the finite principal-open cover $D(h_i)$.
\end{proof}

\begin{theorem}[Principalization with empty boundary]\label{thm:sc0-algebraic-principalization}
\lean{Hironaka.Sequence.BP_isIdealOfSncDivisor,Hironaka.Sequence.BP_center_smooth_hasSncWith,Hironaka.Sequence.BP_isIso_restrict}
\leanok
Let $k$ have characteristic zero, let $X$ be smooth, equidimensional, separated and of finite
type over $k$, and let $I$ be a coherent ideal nonzero at every stalk. There is a finite sequence
of blowups $\pi\colon Y\to X$ with smooth centres having simple normal crossings with the
successive exceptional boundaries, such that $I\mathcal O_Y$ is the ideal of an effective divisor
supported on a simple normal crossing family, and $\pi$ is an isomorphism over $X\setminus V(I)$.
\end{theorem}
\begin{proof}
Apply the principalization construction of Koll\'ar, Theorem~35 of \cite{Kollar07}, to the triple
$(X,I,\varnothing)$. Its order-reduction construction gives the finite sequence and the smooth
normal-crossing centres. At mark one the last marked transform is the unit ideal; the total
transform is therefore the product of the transformed exceptional ideals, which is an effective
simple normal crossing divisor. The isomorphism locus in the general ordered-boundary statement
is $X\setminus(V(I)\cup\operatorname{Sing}E)$. Since the boundary is empty here, this is exactly
$X\setminus V(I)$.
\end{proof}

\begin{lemma}[A finite tree of principalization charts]\label{lem:sc0-bp-chart-tree}
\notready
\uses{thm:sc0-algebraic-principalization,lem:sc0-hironaka-affine-chart,lem:sc0-principal-cover-chart}
For a smooth finite-type domain $B$ over a characteristic-zero field and a nonzero ideal
$I\subseteq B$, the principalization above admits a finite Rees chart tree rooted at $B$ whose
leaves are affine open subschemes of $Y$, with their coordinate rings embedded in $\operatorname{Frac}B$.
\end{lemma}
\begin{proof}
The isomorphism over the nonempty dense set $D(I)$ identifies the function fields. At the first
stage restrict the centre to $\operatorname{Spec}B$ and choose a finite set of generators of its
ideal, which exists because $B$ is Noetherian. This ideal is nonzero: blowing up the zero ideal
would give the empty scheme, contradicting the isomorphism over $D(I)$. The affine charts of this
blowup are exactly the charts in Lemma~\ref{lem:sc0-hironaka-affine-chart}. For every nonzero
denominator, restrict the remaining sequence to that chart and repeat. Restriction preserves its
isomorphism over the generic point, so the centre at each nonempty chart is again nonzero.
Induction on the number of remaining blowups constructs finite subtrees. Each finite generator
family gives finitely many children, so the resulting tree is finite. Any additional finite
principal-open refinements are attached using Lemma~\ref{lem:sc0-principal-cover-chart}.
\end{proof}

\begin{lemma}[Divisors of a product of prime elements]\label{lem:sc0-prime-monomial-factor}
\lean{EnvelopeProperty.Char0.exists_exponents_associated_of_dvd_monomial}
\leanok
Let $R$ be a domain, let $p_1,\ldots,p_s$ be prime elements, and let $e_i\in\mathbf N$.
If $f\mid\prod_i p_i^{e_i}$, then $f$ is associated to $\prod_i p_i^{a_i}$ for suitable
$0\le a_i\le e_i$.
\end{lemma}
\begin{proof}
Induct on $N=\sum_i e_i$. If $N=0$, then $f\mid1$, so $f$ is a unit and all $a_i=0$ work.
Otherwise choose $j$ with $e_j>0$ and write the product as $p_jM$. Choose $q$ with $fq=p_jM$.
Primality gives $p_j\mid f$ or $p_j\mid q$. In the first case write $f=p_jf'$ and cancel the
nonzero $p_j$ to obtain $f'q=M$; apply induction to $f'$ and increase its $j$-th exponent by one.
In the second case write $q=p_jq'$ and cancel to obtain $fq'=M$; apply induction to $f$ and leave
its exponents unchanged. In both cases the resulting exponents satisfy the required bounds.
\end{proof}

\begin{lemma}[Units away from a prime boundary]\label{lem:sc0-unit-away-monomial}
\lean{EnvelopeProperty.Char0.exists_exponents_associated_of_isUnit_away_product,EnvelopeProperty.Char0.exists_unit_mul_monomial_of_isUnit_away_product}
\leanok
\uses{lem:sc0-prime-monomial-factor}
Let $R$ be a domain and let $p_1,\ldots,p_s$ be prime elements. If the image of $f\in R$
is a unit in $R[1/(\prod_i p_i)]$, then $f=u\prod_i p_i^{a_i}$ for some unit $u\in R^\times$
and nonnegative integers $a_i$.
\end{lemma}
\begin{proof}
The localization criterion for units gives $f\mid(\prod_i p_i)^N$ for some $N$.
Distribute the power over the product and apply Lemma~\ref{lem:sc0-prime-monomial-factor}
with all exponents equal to $N$. Associated elements differ by a unit, giving the stated form.
For an empty boundary the product is one, and the same argument says that $f$ is a unit.
\end{proof}

\begin{lemma}[Adapted affine \'etale coordinates]\label{lem:sc0-snc-etale-charts}
\lean{EnvelopeProperty.Char0.exists_snc_etale_chart_at_closed,EnvelopeProperty.Char0.exists_finite_snc_etale_cover}
\leanok
Let $Y$ be smooth of dimension $d$ over a characteristic-zero field and let $D$ be a finite
simple normal crossing divisor family. There is a cover of $Y$ by affine neighbourhoods with
\'etale coordinates $u_1,\ldots,u_d$ on which each member of $D$ is either absent or has ideal
$(u_i)$, with distinct components assigned distinct indices. A quasi-compact $Y$ admits a finite
such cover.
\end{lemma}
\begin{proof}
At a closed point the ideals of the boundary components are generated by distinct members
of a regular system of $d$ parameters. Their classes form a basis of the cotangent space.
The residue field is finite and separable over the characteristic-zero base, so the conormal
sequence identifies this cotangent space with the fibre of $\Omega_{Y/k}$. Represent the
finitely many parameter germs on a common affine neighbourhood. The resulting map to
$\mathbf A_k^d$ is \'etale where the differential determinant is invertible, an open neighbourhood
of the closed point.
Equality of the finitely generated component ideals with the corresponding principal ideals
also holds on a neighbourhood; remove every component not passing through the point and intersect
these finitely many neighbourhoods. Refine the resulting open set by a principal affine open.
The scheme is Jacobson because it is locally of finite type over a field. Every point specializes
to a closed point, and an open neighbourhood of that closed point contains its generalizations.
Consequently these neighbourhoods cover $Y$.
Quasi-compactness selects finitely many members of this cover.
\end{proof}

\begin{lemma}[Stalks of factors of an SNC product]\label{lem:sc0-snc-factor-stalks}
\lean{EnvelopeProperty.Char0.prod_pow_stalkIdeal_eq_span,EnvelopeProperty.Char0.exists_stalk_monomial_factor}
\leanok
\uses{lem:sc0-prime-monomial-factor}
Let $A,B$ be ideal sheaves on a scheme smooth over a field, and suppose
$AB=\prod_i F_i^{e_i}$ for a simple normal crossing family $F$. At every point $x$, choose
regular parameters $z$ cutting out the members of $F$ through $x$. Then $A_x$ is generated
by $\prod_{i:x\in F_i} z_{c(i)}^{a_i}$ with $0\le a_i\le e_i$.
\end{lemma}
\begin{proof}
The SNC product is invertible, so both factors are invertible. Hence $A_x=(g)$ for a nonzero
generator $g$. The members of $F$ not through $x$ have unit stalk ideal; the remaining members
have ideals $(z_{c(i)})$. Consequently $(AB)_x$ is generated by
$m=\prod_{i:x\in F_i}z_{c(i)}^{e_i}$. Since $(AB)_x\subseteq A_x$, the generator $g$ divides $m$.
The stalk is a regular local domain, and each boundary parameter is prime. Apply
Lemma~\ref{lem:sc0-prime-monomial-factor}. Associated generators have the same principal ideal.
\end{proof}

\begin{lemma}[Monomial ideals from factors of the total transform]\label{lem:sc0-snc-factors}
\notready
\uses{lem:sc0-snc-factor-stalks,lem:sc0-snc-etale-charts}
If a product of finitely many nonzero ideal sheaves on a smooth integral scheme is the ideal of
an effective simple normal crossing divisor, then each factor is, on a sufficiently small adapted
affine neighbourhood, generated by a monomial in the same boundary coordinates.
\end{lemma}
\begin{proof}
For one factor let $B$ be the product of the others. Apply
Lemma~\ref{lem:sc0-snc-factor-stalks} to obtain its monomial stalk ideal in the same boundary
parameters. Repeat for the finitely many factors. Choose adapted \'etale coordinates as in
Lemma~\ref{lem:sc0-snc-etale-charts}. Their assigned boundary functions generate the same
component stalk ideals, so replacing the original parameters preserves each monomial ideal.
Equalities of finitely
generated ideals and invertibility of the finitely many units spread to one affine principal
neighbourhood. Multiplying by a unit does not change a principal ideal.
\end{proof}

\begin{lemma}[The relative Jacobian ideal is an invertible image ideal]\label{lem:sc0-jacobian-line}
\notready
Let $\pi\colon Y\to X$ be a dominant generically isomorphic morphism of smooth integral
$d$-dimensional varieties. In \'etale coordinates $u$ on $Y$, the ideal generated by
$\det(\partial(\pi^*w_a)/\partial u_b)$ for regular functions $w_1,\ldots,w_d$ on $X$ is
locally a nonzero principal ideal. It is the unit ideal over the isomorphism locus of $\pi$.
\end{lemma}
\begin{proof}
Refine the image neighbourhood in $X$ by \'etale coordinates $z$. The differential of $\pi$
between the free differential modules has a square matrix $A$ in the $dz$ and $du$ bases.
The determinant for a family $w$ is $\det(B_wA)=\det(B_w)\det A$, so every generator lies in
$(\det A)$. Conversely the $dz_i$ are linear combinations of the differentials of functions
from the original affine ring, after multiplying denominators which are units on this
neighbourhood. Multilinearity of the determinant expresses $\det A$ as a linear combination
of the displayed generators. Thus the ideal is $(\det A)$. Dominance and the generic isomorphism
make $A$ invertible at the function field, so $\det A\ne0$. If $\pi$ is an isomorphism, its inverse
differential makes $A$ invertible already in the coordinate ring, so the determinant ideal is the
unit ideal. This description is independent of the chosen bases because a change of basis has
unit determinant.
\end{proof}

\begin{lemma}[The relative Jacobian is a boundary monomial]\label{lem:sc0-jacobian-snc}
\notready
\uses{lem:sc0-jacobian-line,lem:sc0-unit-away-monomial,lem:sc0-snc-etale-charts}
If $\pi$ in the preceding lemma is an isomorphism away from a simple normal crossing divisor
$D$ on $Y$, then on a sufficiently small adapted affine neighbourhood its relative Jacobian
ideal is generated by a monomial in the boundary coordinates.
\end{lemma}
\begin{proof}
At a stalk choose adapted boundary parameters $u_1,\ldots,u_s$ and a nonzero generator $J$ of
the relative Jacobian ideal. Let $q=\prod_i u_i$. The generator is a unit in $R[1/q]$, because
this open set avoids the boundary and $\pi$ is an isomorphism there. Each $u_i$ is prime, so
Lemma~\ref{lem:sc0-unit-away-monomial} identifies $J$ with a
boundary monomial up to a unit. Spread this finite equality and the inverse of its unit to a
principal affine neighbourhood. The empty-boundary case has $q=1$ and gives the unit ideal.
\end{proof}

\begin{theorem}[Finite log-chart trees in characteristic zero]\label{thm:sc0-hironaka-log-tree}
\lean{EnvelopeProperty.Char0.exists_smooth_family_rees_tree}
\leanok
\uses{thm:sc0-smooth-family-rees-tree}
For a smooth finite-type domain $B$ over a characteristic-zero field and a finite family of
nonzero ideals $\mathfrak a_j$, there is a finite Rees chart tree over $B$ whose leaves have
\'etale coordinates in which each $\mathfrak a_j$ and the relative Jacobian ideal are monomial.
\end{theorem}
\begin{proof}
Regard $B$ as a subalgebra of its fraction field. Choose finitely many principal
source opens, generating the unit ideal, on which the smooth source has
\'etale coordinates. These opens form the first, unit-ideal Rees node. Each
localized family factor is nonzero, since the localization maps embed in the
same fraction field. On each source chart principalize the product of all
family factors by the actual characteristic-zero blowup sequence. Its final
scheme is smooth, the total transform is a simple normal crossing divisor,
and the morphism is an isomorphism off the original product support.

The adapted simultaneous factor and Jacobian neighbourhoods cover this final
scheme. Refine its finite Rees chart tree subordinate to those neighbourhoods.
On every resulting literal leaf, transport the boundary-ideal identities to
the leaf coordinates. Every family factor is then a monomial in these
coordinates, and the determinant formula identifies the genuine relative
Jacobian ideal generated by all source functions with a boundary monomial.
Attach these child trees to the initial source-cover node. Localization of
the source preserves both ideal extension and the genuine relative Jacobian
ideal, so these formulas refer to the original $B$ and its original family.
\end{proof}
